\chapter{Introducción}
\label{cap:introduccion}

\section{Justificación y motivación}

El vídeo se ha consolidado como uno de los principales soportes de aprendizaje. Plataformas
como YouTube concentran una enorme cantidad de clases, charlas y tutoriales, pero el formato
audiovisual presenta una limitación clara para el estudio: es lineal y poco consultable. Para
localizar un concepto concreto dentro de una grabación de una hora, el estudiante debe recordar
aproximadamente en qué momento se trató y desplazarse manualmente por la barra de
reproducción. A diferencia de un texto, un vídeo no se puede buscar, ni resumir, ni hojear con
facilidad.

Existen herramientas parciales —generación de subtítulos, capítulos manuales o resúmenes
automáticos—, pero rara vez se integran en un único flujo pensado para el estudio. Lo habitual
es que el alumno combine varias aplicaciones y siga sin poder \textit{preguntar} directamente
al contenido del vídeo.

ComprendiA nace para cubrir ese hueco: convertir cualquier vídeo de YouTube en material de
estudio estructurado y consultable. La aplicación transcribe la grabación, la analiza para
extraer capítulos y conceptos clave, y ofrece un asistente conversacional capaz de responder
preguntas sobre el contenido citando el momento exacto del vídeo. Además, organiza las clases
en asignaturas y reduce el trabajo administrativo sugiriendo automáticamente a qué asignatura
pertenece cada nueva grabación.

\section{Objetivos}

\subsection{Objetivo general}

Diseñar e implementar una aplicación web que permita transformar vídeos de YouTube en material
de estudio estructurado, ofreciendo transcripción, análisis automático del contenido y consulta
mediante un asistente conversacional basado en recuperación semántica.

\subsection{Objetivos específicos}

\begin{itemize}
  \item Obtener la transcripción de un vídeo de YouTube de forma eficiente, priorizando los
        subtítulos existentes y recurriendo a un modelo de reconocimiento del habla cuando sea
        necesario.
  \item Indexar el contenido mediante \textit{embeddings} para permitir búsqueda semántica por
        fragmentos.
  \item Generar automáticamente capítulos navegables y conceptos clave con sus marcas de tiempo.
  \item Implementar un asistente conversacional (RAG) que responda preguntas sobre el vídeo con
        memoria corta del diálogo y referencias al momento concreto de la grabación.
  \item Organizar las clases en asignaturas y sugerir automáticamente la asignatura de cada
        vídeo nuevo mediante criterios de canal y de similitud semántica.
  \item Construir una interfaz web clara que integre reproductor, capítulos, conceptos y chat.
\end{itemize}

\section{Alcance del proyecto}

El proyecto abarca el desarrollo completo de la aplicación, tanto del \textit{backend} (API,
procesamiento y persistencia) como del \textit{frontend} (interfaz de usuario). Se contempla el
procesamiento de vídeos públicos de YouTube y la gestión de clases y asignaturas por parte de un
único perfil de usuario.

Quedan fuera del alcance de esta primera versión la gestión multiusuario con autenticación, el
procesamiento de fuentes de vídeo distintas a YouTube y el despliegue en un entorno de
producción con alta disponibilidad. Estos puntos se recogen como líneas futuras en el
capítulo~\ref{cap:conclusiones}.

\section{Estructura de la memoria}

El resto de la memoria se organiza como sigue. El capítulo~\ref{cap:estado-arte} revisa el
estado del arte de las tecnologías implicadas (transcripción automática, \textit{embeddings},
búsqueda semántica y modelos de lenguaje). El capítulo~\ref{cap:requisitos} recoge la
especificación de requisitos del sistema. El capítulo~\ref{cap:desarrollo} describe la
metodología, la arquitectura y la implementación. El capítulo~\ref{cap:pruebas} presenta las
pruebas realizadas y sus resultados. Por último, el capítulo~\ref{cap:conclusiones} expone las
conclusiones y las líneas de trabajo futuras.
